\section{Human Evaluation}
\label{sec:human}


To verify that our automatic evaluation correlates with human judgement, 
we conduct human evaluation on selected models %
and request workers 
to judge model generations on three dimensions similar to \citet{liu2023evaluating}---(1) utility: a 1-to-5 score indicating whether the generation helps answer the question; %
(2) citation recall: the annotator is given a sentence and all passages that the sentence cited, and is asked to judge whether the passages fully support the sentence; %
(3) citation precision: given a sentence and one of its citations, the annotator is  asked to judge whether the citation ``fully supports'', ``partially supports'', or ``does not support'' the sentence. 
Each citation gets a precision score 1 if the output sentence has a citation recall of 1 and this citation at least ``partially supports'' it.
See Appendix~\ref{app_sec:human} for more details.


\begin{table}[t]
    \centering
    \resizebox{0.98\linewidth}{!}{
    \begin{tabular}{lcccc}
        \toprule
            & \multicolumn{2}{c}{\tf{Human scores}} & \multicolumn{2}{c}{\tf{\citeeval{} scores}} \\
             \cmidrule(lr){2-3} \cmidrule(lr){4-5}
            & Rec.& Prec. & Rec. & Prec. \\
            \midrule
            ChatGPT \vani & 74.7 & 76.6 & 75.3 & 74.4 \\
            ~~w/ \rerank & \textbf{79.3} & \textbf{81.9} & \textbf{83.9} & \textbf{80.8} \\
            Vicuna-13B \vani & 51.6 & 51.5 & 50.3 & 50.1 \\
            \bottomrule
            \end{tabular}
    }
    \caption{
        Human citation quality evaluation vs. \citeeval{} citation quality evaluation on ASQA. 
    }
    \vspace{-5pt}
    \label{tab:human-asqa-all}
\end{table}

\begin{table}[t]
    \centering
    \resizebox{0.98\linewidth}{!}{
    \begin{tabular}{lcccc}
        \toprule
            & \multicolumn{2}{c}{\tf{Human scores}} & \multicolumn{2}{c}{\tf{\citeeval{} scores}} \\
             \cmidrule(lr){2-3} \cmidrule(lr){4-5}
             & Rec. & Prec. & Rec. & Prec. \\
            \midrule
            ChatGPT \vani & 50.8 & 52.4 & 52.8 & 50.4 \\
            ~~w/ \rerank & \textbf{59.7} & \textbf{60.6} & \textbf{63.0} & \textbf{60.6} \\
            Vicuna-13B \vani & 13.4 & 19.2 & 13.6 & 18.1 \\
            \bottomrule
            \end{tabular}
    }
    \caption{
        Human citation quality evaluation vs. \citeeval{} citation quality evaluation on ELI5. 
    }
    \vspace{-10pt}
    \label{tab:human-eli5-all}
\end{table}


\paragraph{Model outputs score high utility.} 
The utility scores do not differ significantly between models, ranging $3.7$-$3.9$ for ASQA and $3.5$-$3.6$  for ELI5. 
Upon inspection, all tested models are mostly able to output fluent answers that are related to the question, %
despite  differences in factual correctness. 





\paragraph{Our automatic evaluation of citation quality strongly correlates with human judgements.}
As shown in Table~\ref{tab:human-asqa-all} (ASQA) and Table~\ref{tab:human-eli5-all} (ELI5),
the relative rankings induced by human and our automatic metrics are consistent.
The absolute citation scores from human and \citeeval{} are very close except for \rerank{} (which uses the automated citation recall for reranking).
This suggests that an improvement on \citeeval{} citation metrics translates to improvement on human preferences. %
Furthermore, the Cohen's kappa coefficient between 
human and \citeeval{} suggests substantial agreement for citation recall ($0.698$) and moderate agreement for citation precision ($0.525$).
We also show in \S\ref{sec_more_human} that our automatic evaluation achieves high accuracy when treating human annotations as gold labels
($85.1\%$ for citation recall and $77.6\%$ for citation precision).










